\chapter{Catalogue des scénarios}
\label{ann:scenarios}
% BUDGET : 2 pages. Version 2026-09-02 : l'annexe ne contient plus que des tableaux. Les totaux
% de reference, les remarques de lecture et les regles de balayage, jusqu'ici en prose, sont
% devenus trois tables.

\section{Les totaux dont les seuils sont tirés}

\begin{table}[H]
  \centering
  \caption{Totaux mesurés sur la calibration retenue. Chaque seuil des catalogues qui suivent
  est une fraction déclarée de l'une de ces valeurs, et non un chiffre choisi arbitrairement.}
  \label{tab:ann-totaux}
  \small
  \begin{tabularx}{\textwidth}{@{}>{\raggedright\arraybackslash}Xr@{}}
    \toprule
    Marge brute du territoire & \SI{\calibRetenuMarge}{\euros} \\
    Subventions versées & \SI{\calibRetenuSubvention}{\euros} \\
    Emploi & \SI{\calibRetenuEtp}{\etp} \\
    Azote épandu & \SI{\calibRetenuAzote}{\kilogram} \\
    \gls{ift} total & \num{\calibRetenuIft} \\
    Eau prélevée, parcelles irrigables seules & \SI{\calibRetenuEau}{\metre\cubed} \\
    \bottomrule
  \end{tabularx}
\end{table}

\section{Les dix politiques}

\begin{table}[H]
  \centering
  \caption{Les dix politiques du catalogue et leurs leviers, du pôle accélérationniste à la
  rupture agroécologique. Aides désigne un multiplicateur appliqué aux subventions d'un groupe
  de cultures, \emph{slack} le desserrement du plafond de main d'\oe uvre, inertie la part de
  surface qui doit rester dans son groupe de 2017.}
  \label{tab:ann-politiques}
  \scriptsize
  \setlength{\tabcolsep}{4pt}
  \begin{tabularx}{\textwidth}{@{}p{27mm}X@{}}
    \toprule
    \textbf{Politique} & \textbf{Leviers} \\
    \midrule
    P1 dérégulation totale & aides supprimées, marge brute pure, quotas, planchers, rotations
      et interdictions désactivés, slack \num{3.0}. Borne supérieure, et mesure de ce que
      l'ensemble des règles coûte \\
    P2 accélérationnisme industriel & aides canne $\times$\num{1.5}, banane $\times$\num{1.4},
      quotas d'export levés, jachère de moitié, part mécanisée \SI{60}{\percent}, slack
      \num{0.85} \\
    P3 intensification modérée & aides $\times$\num{1.20} et $\times$\num{1.15}, quotas
      $+\SI{30}{\percent}$, enveloppe plafonnée, inertie \SI{40}{\percent} \\
    P4 statu quo & règles inchangées, enveloppe, \gls{ift} et azote gelés au niveau constaté,
      plancher d'emploi, inertie \SI{70}{\percent}. Point de comparaison des autres \\
    P5 austérité budgétaire & enveloppe $-\SI{40}{\percent}$, coupe supplémentaire sur la
      canne, aucune compensation. Mesure ce que la géographie doit aux aides plutôt qu'à
      l'agronomie \\
    P6 verdissement incitatif & la carotte seule. Surprime au maraîchage biologique,
      réouverture des systèmes qui le permettent, aucune interdiction nouvelle, slack
      \num{1.20} \\
    P7 Écophyto réglementaire & le bâton seul. \gls{ift} $-\SI{30}{\percent}$, azote
      $-\SI{20}{\percent}$, nitrates \SI{170}{\kilogram\per\ha} par ferme, prairie interdite
      sur sols chlordéconés. Le pendant exact de P6 \\
    P8 transition agroécologique & \gls{ift} $-\SI{40}{\percent}$, azote
      $-\SI{30}{\percent}$, eau $-\SI{20}{\percent}$, nitrates \SI{150}{\kilogram\per\ha},
      aides de l'export vers le biologique, emploi $\geq\SI{4200}{\etp}$, slack \num{1.50},
      inertie \SI{40}{\percent} \\
    P9 souveraineté alimentaire & hors axe. Planchers vivriers réactivés, quota sucrier
      abaissé, aides vers le local, plancher arboricole, slack \num{1.60}, inertie
      \SI{30}{\percent} \\
    P10 bifurcation agroécologique & symétrique de P1. \gls{ift} $-\SI{50}{\percent}$, azote
      $-\SI{40}{\percent}$, eau $-\SI{30}{\percent}$, \gls{ges} $-\SI{25}{\percent}$, export
      divisé par deux, aucune inertie. Aucune solution entière trouvée
      (figure~\ref{fig:spectre}) \\
    \bottomrule
  \end{tabularx}
\end{table}

\section{Les douze forçages}

\begin{table}[H]
  \centering
  \caption{Les douze forçages du catalogue. Un forçage multiplie des prix, des rendements, des
  variances ou des coûts d'intrants et ne modifie aucune règle, à l'exception de F11 qui retire
  de la terre. F0 est le forçage nul, F10 la borne haute symétrique de F9.}
  \label{tab:ann-forcages}
  \scriptsize
  \setlength{\tabcolsep}{4pt}
  \begin{tabularx}{\textwidth}{@{}p{31mm}X@{}}
    \toprule
    \textbf{Forçage} & \textbf{Ce qu'il impose} \\
    \midrule
    F0 nominal & aucun forçage. Conditions 2017 \\
    F1 cyclone majeur & rendements $\times$\num{0.50} à $\times$\num{0.75} selon le port,
      variance $\times$\num{1.30}. Les pérennes hautes paient le plus \\
    F2 sécheresse sévère & rendements $-\num{25}$ à $-\SI{35}{\percent}$, variance
      $\times$\num{1.40}, prélèvement d'eau $-\SI{30}{\percent}$. Éprouve la dépendance à
      l'eau \\
    F3 dérive climatique & rendements $-\SI{12}{\percent}$, variance $+\SI{60}{\percent}$.
      Isole le canal du risque \\
    F4 choc sanitaire & cercosporiose noire. Banane d'export $-\SI{50}{\percent}$ de rendement
      et $+\SI{25}{\percent}$ de coût de protection. Choc de filière \\
    F5 effondrement des prix d'export & banane $-\SI{35}{\percent}$, canne
      $-\SI{25}{\percent}$. Les deux piliers cèdent ensemble \\
    F6 choc sur les intrants & coûts variables $+\SI{45}{\percent}$ partout. Sépare le pôle
      accélérationniste du pôle agroécologique \\
    F7 inflation alimentaire & vivrier local $+\SI{40}{\percent}$, export inchangé. Favorable
      à la relocalisation \\
    F8 retrait du \gls{posei} & aides à \SI{40}{\percent} de leur niveau. Choc purement
      budgétaire \\
    F9 crise systémique & rendements $-\num{20}$, prix d'export $-\num{25}$, coûts $+\num{30}$,
      subventions $-\SI{40}{\percent}$, variance $+\SI{50}{\percent}$. Le pire cas plausible \\
    F10 conjoncture favorable & rendements $+\num{10}$, prix $+\num{15}$, coûts $-\num{10}$,
      variance $-\SI{20}{\percent}$. Le témoin par le haut \\
    F11 artificialisation & \SI{8}{\percent} de la \gls{sau} sortent de la production, en
      priorité à basse altitude. Le seul forçage qui retire de la terre \\
    \bottomrule
  \end{tabularx}
\end{table}

\section{Conventions de lecture et règles de campagne}

\begin{table}[H]
  \centering
  \caption{Les cinq conventions qui rendent les deux catalogues lisibles, et ce que chacune
  implique pour la comparaison des résultats.}
  \label{tab:ann-conventions}
  \scriptsize
  \setlength{\tabcolsep}{4pt}
  \begin{tabularx}{\textwidth}{@{}>{\raggedright\arraybackslash}p{34mm}X@{}}
    \toprule
    \textbf{Convention} & \textbf{Conséquence} \\
    \midrule
    P1 change d'objectif & elle ne se compare aux autres que sur les grandeurs physiques et
      économiques, jamais sur la colonne d'objectif \\
    Aucune politique ne rouvre les variantes symétriques & leur réouverture coûterait
      \num{548000} binaires de pure symétrie et rendrait toute part biologique inopérante
      (§~\ref{sec:tractabilite}). Écrire la politique agroécologique comme une part minimale de
      cultures biologiques en était la formulation naturelle, et c'est le seul cas du mémoire
      où une formulation est écartée parce qu'elle serait satisfaite trop facilement \\
    Le plafond de main d'\oe uvre est la première cause d'infaisabilité & il accorde
      \SI{3891}{\etp} au desserrement nominal, donc tout plancher d'emploi supérieur exige de
      relever le desserrement dans le même scénario. Les planchers du catalogue sont posés vers
      \SI{86}{\percent} à \SI{88}{\percent} d'un maximum mesuré en maximisant le travail sous
      les autres contraintes de la politique, pour laisser la place au saut d'intégralité que
      la relaxation ignore \\
    Les forçages ne sont pas moyennés & le résultat d'une politique est un vecteur et non un
      scalaire. L'infaisabilité y compte comme le pire résultat et non comme une donnée
      manquante \\
    F0 est le forçage nul & une politique seule et la même politique croisée avec F0 sont le
      même problème et doivent rendre le même chiffre. Ce contrôle a détecté une erreur
      d'ancrage en cours de campagne, les deux runs différant de \SI{\prosEcartFzero}{\euros},
      soit \SI{\prosEcartFzeroPct}{\percent}. Il ne coûte rien et doit être refait à chaque
      campagne \\
    \bottomrule
  \end{tabularx}
\end{table}

\begin{table}[H]
  \centering
  \caption{Règles de construction des balayages et contrôles appliqués avant qu'une campagne ne
  tourne. Deux fronts sont déclarés, l'un sur le plafond d'azote et l'autre sur l'enveloppe de
  subventions.}
  \label{tab:ann-balayages}
  \scriptsize
  \setlength{\tabcolsep}{4pt}
  \begin{tabularx}{\textwidth}{@{}>{\raggedright\arraybackslash}p{34mm}X@{}}
    \toprule
    \textbf{Règle ou contrôle} & \textbf{Contenu} \\
    \midrule
    Un front n'est pas une cellule & c'est une résolution par point. Il n'est donc jamais
      croisé avec les forçages de sa cellule, et un front forcé doit être nommé explicitement \\
    Un front se paramètre sur sa politique hôte & et non sur le réalisé de la configuration de
      référence. Verdict renversé en cours de campagne, dont l'oubli produit un front dont
      plusieurs points ne contraignent rien \\
    Les points sont réordonnés du plus serré au plus lâche & à l'inverse de la lecture du
      fichier, pour que chaque solution serve de graine admissible à la suivante. Le sens
      s'inverse pour un plancher \\
    Le tri s'abstient dès qu'il ne peut pas conclure & s'abstenir ne coûte que du temps de
      calcul, la correction reposant sur l'audit de chaque graine \\
    Contrôle de faisabilité en relaxation linéaire & deux minutes contre une heure. Infaisable
      en relaxation implique infaisable en entiers, sans que la réciproque tienne. C'est donc
      un filtre et non une garantie. Il a écarté une erreur de conception avant qu'aucun lot ne
      tourne, deux blocs corrects séparément ne composant pas nécessairement un scénario
      faisable \\
    Détection des groupes de cultures symétriques & le même contrôle les signale, et c'est
      ainsi que les vingt-cinq variantes ont été trouvées \\
    Piège du poids de parcelle & un seuil relevé sur un run et reposé en configuration sans son
      poids de parcelle est faux d'un facteur \num{1.6}. C'est le cas du plafond d'eau, qui ne
      porte que sur les parcelles irrigables. Ce piège ne laisse aucune trace dans les
      sorties \\
    \bottomrule
  \end{tabularx}
\end{table}
